% Part: computability
% Chapter: recursive-functions
% Section: halting-problem

\documentclass[../../../include/open-logic-section]{subfiles}

\begin{document}

\olfileid[pt]{cmp}{rec}{hlt}
\olsection{O problema da parada}

O \emph{problema da parada}, em geral, é o problema de decidir, dada a
especificação~$e$ (por exemplo, um programa) de uma função computável e
um número~$n$, se a computação da função na entrada~$n$ para, isto é,
se produz um resultado. Alan Turing demonstrou, em um resultado célebre,
que esse próprio problema não pode ser resolvido por uma função
computável, isto é, a função
\[
h(e, n) =
\begin{cases}
1 & \text{se a computação $e$ para na entrada $n$}\\
0 & \text{caso contrário,}
\end{cases}
\]
não é computável.

No contexto das funções recursivas parciais, o papel da especificação
de um programa pode ser desempenhado pelo índice~$e$ dado no teorema
da forma normal de Kleene. Se $f$ é uma função recursiva parcial,
qualquer $e$~para o qual vale a equação do teorema da forma normal é
um índice de~$f$. Dado um número~$e$, o teorema da forma normal afirma
que
\[
\cfind{e}(x) \simeq U(\mu s \; T(e, x, s))
\]
é recursiva parcial e, para toda função recursiva parcial $f\colon \Nat
\to \Nat$, existe um $e \in \Nat$ tal que $\cfind{e}(x) \simeq
f(x)$ para todo~$x \in \Nat$. Na verdade, para cada função~$f$ assim
há não apenas um, mas infinitos~$e$ com essa propriedade. A
\emph{função de parada}~$h$ é definida por
\[
h(e, x) =
\begin{cases}
1 & \text{se $\cfind{e}(x) \fdefined$}\\
0 & \text{caso contrário.}
\end{cases}
\]
Observe que $h(e, x) = 0$ se $\cfind{e}(x) \fundefined$, mas também
quando~$e$ não é índice de nenhuma função recursiva parcial.

\begin{thm}
\ollabel{thm:halting-problem}
A função de parada~$h$ não é recursiva parcial.
\end{thm}

\begin{proof}
Se $h$
fosse recursiva parcial, poderíamos definir
\[
d(y) =
\begin{cases}
1 & \text{se $h(y, y) = 0$}\\
\umin{x}{x \neq x} & \text{caso contrário.}
\end{cases}
\]
Como nenhum número $x$ satisfaz $x \neq x$, não existe
$\umin{x}{x \neq x}$; portanto, $d(y) \fundefined$ se, e somente se,
$h(y,y) \neq 0$. Dessa definição segue que
\begin{enumerate}
\item $d(y) \fdefined$ se, e somente se, $\cfind{y}(y) \fundefined$
  ou $y$ não é índice de uma função recursiva parcial.
\item $d(y) \fundefined$ se, e somente se, $\cfind{y}(y) \fdefined$.
\end{enumerate}
Se $h$ fosse recursiva parcial, então $d$ também seria recursiva
parcial. Assim, pelo teorema da forma normal de Kleene, ela teria um
índice~$e_d$. Considere o valor de $h(e_d, e_d)$. Há dois casos
possíveis, $0$ e~$1$.
\begin{enumerate}
\item Se $h(e_d, e_d) = 1$, então $\cfind{e_d}(e_d) \fdefined$. Mas
  $\cfind{e_d} \simeq d$, e $d(e_d)$ está definida se, e somente se,
  $h(e_d, e_d) = 0$. Logo, $h(e_d, e_d) \neq 1$.
\item Se $h(e_d, e_d) = 0$, então ou $e_d$ não é índice de uma função
  recursiva parcial, ou ele é e $\cfind{e_d}(e_d) \fundefined$. Mas,
  novamente, $\cfind{e_d} \simeq d$, e $d(e_d)$ é indefinida se, e
  somente se, $\cfind{e_d}(e_d) \fdefined$.
\end{enumerate}
O resultado é que $e_d$ não pode, afinal, ser o índice de uma função
recursiva parcial. Mas, se $h$ fosse recursiva parcial, $d$ também
seria, e nossa definição de $e_d$ como um índice dela seria
admissível. Devemos concluir que $h$ não pode ser recursiva parcial.
\end{proof}

\end{document}

